%!TEX TS-program = xelatex
\documentclass[10pt,oneside]{article}
\usepackage[fontsize=9pt]{scrextend}

\usepackage[english]{babel}

\usepackage{amsmath,amssymb,amsfonts}
\usepackage[utf8]{inputenc}
\usepackage[T1]{fontenc}
\usepackage{stix2}
\usepackage[scaled]{helvet}
\usepackage[scaled]{inconsolata}

\usepackage{lastpage}

\usepackage{setspace}

\usepackage{ccicons}

\usepackage[hang,flushmargin]{footmisc}

\usepackage{geometry}

\setlength{\parindent}{0pt}
\setlength{\parskip}{6pt plus 2pt minus 1pt}

\usepackage{fancyhdr}
\renewcommand{\headrulewidth}{0pt}\providecommand{\tightlist}{%
  \setlength{\itemsep}{0pt}\setlength{\parskip}{0pt}}

\makeatletter
\newcounter{tableno}
\newenvironment{tablenos:no-prefix-table-caption}{
  \caption@ifcompatibility{}{
    \let\oldthetable\thetable
    \let\oldtheHtable\theHtable
    \renewcommand{\thetable}{tableno:\thetableno}
    \renewcommand{\theHtable}{tableno:\thetableno}
    \stepcounter{tableno}
    \captionsetup{labelformat=empty}
  }
}{
  \caption@ifcompatibility{}{
    \captionsetup{labelformat=default}
    \let\thetable\oldthetable
    \let\theHtable\oldtheHtable
    \addtocounter{table}{-1}
  }
}
\makeatother

\usepackage{array}
\newcommand{\PreserveBackslash}[1]{\let\temp=\\#1\let\\=\temp}
\let\PBS=\PreserveBackslash

\usepackage[breaklinks=true]{hyperref}
\hypersetup{colorlinks,%
citecolor=blue,%
filecolor=blue,%
linkcolor=blue,%
urlcolor=blue}
\usepackage{url}

\usepackage{caption}
\setcounter{secnumdepth}{0}
\usepackage{cleveref}

\usepackage{graphicx}
\makeatletter
\def\maxwidth{\ifdim\Gin@nat@width>\linewidth\linewidth
\else\Gin@nat@width\fi}
\makeatother
\let\Oldincludegraphics\includegraphics
\renewcommand{\includegraphics}[1]{\Oldincludegraphics[width=\maxwidth]{#1}}

\usepackage{longtable}
\usepackage{booktabs}

\usepackage{color}
\usepackage{fancyvrb}
\newcommand{\VerbBar}{|}
\newcommand{\VERB}{\Verb[commandchars=\\\{\}]}
\DefineVerbatimEnvironment{Highlighting}{Verbatim}{commandchars=\\\{\}}
% Add ',fontsize=\small' for more characters per line
\usepackage{framed}
\definecolor{shadecolor}{RGB}{248,248,248}
\newenvironment{Shaded}{\begin{snugshade}}{\end{snugshade}}
\newcommand{\KeywordTok}[1]{\textcolor[rgb]{0.13,0.29,0.53}{\textbf{#1}}}
\newcommand{\DataTypeTok}[1]{\textcolor[rgb]{0.13,0.29,0.53}{#1}}
\newcommand{\DecValTok}[1]{\textcolor[rgb]{0.00,0.00,0.81}{#1}}
\newcommand{\BaseNTok}[1]{\textcolor[rgb]{0.00,0.00,0.81}{#1}}
\newcommand{\FloatTok}[1]{\textcolor[rgb]{0.00,0.00,0.81}{#1}}
\newcommand{\ConstantTok}[1]{\textcolor[rgb]{0.00,0.00,0.00}{#1}}
\newcommand{\CharTok}[1]{\textcolor[rgb]{0.31,0.60,0.02}{#1}}
\newcommand{\SpecialCharTok}[1]{\textcolor[rgb]{0.00,0.00,0.00}{#1}}
\newcommand{\StringTok}[1]{\textcolor[rgb]{0.31,0.60,0.02}{#1}}
\newcommand{\VerbatimStringTok}[1]{\textcolor[rgb]{0.31,0.60,0.02}{#1}}
\newcommand{\SpecialStringTok}[1]{\textcolor[rgb]{0.31,0.60,0.02}{#1}}
\newcommand{\ImportTok}[1]{#1}
\newcommand{\CommentTok}[1]{\textcolor[rgb]{0.56,0.35,0.01}{\textit{#1}}}
\newcommand{\DocumentationTok}[1]{\textcolor[rgb]{0.56,0.35,0.01}{\textbf{\textit{#1}}}}
\newcommand{\AnnotationTok}[1]{\textcolor[rgb]{0.56,0.35,0.01}{\textbf{\textit{#1}}}}
\newcommand{\CommentVarTok}[1]{\textcolor[rgb]{0.56,0.35,0.01}{\textbf{\textit{#1}}}}
\newcommand{\OtherTok}[1]{\textcolor[rgb]{0.56,0.35,0.01}{#1}}
\newcommand{\FunctionTok}[1]{\textcolor[rgb]{0.00,0.00,0.00}{#1}}
\newcommand{\VariableTok}[1]{\textcolor[rgb]{0.00,0.00,0.00}{#1}}
\newcommand{\ControlFlowTok}[1]{\textcolor[rgb]{0.13,0.29,0.53}{\textbf{#1}}}
\newcommand{\OperatorTok}[1]{\textcolor[rgb]{0.81,0.36,0.00}{\textbf{#1}}}
\newcommand{\BuiltInTok}[1]{#1}
\newcommand{\ExtensionTok}[1]{#1}
\newcommand{\PreprocessorTok}[1]{\textcolor[rgb]{0.56,0.35,0.01}{\textit{#1}}}
\newcommand{\AttributeTok}[1]{\textcolor[rgb]{0.77,0.63,0.00}{#1}}
\newcommand{\RegionMarkerTok}[1]{#1}
\newcommand{\InformationTok}[1]{\textcolor[rgb]{0.56,0.35,0.01}{\textbf{\textit{#1}}}}
\newcommand{\WarningTok}[1]{\textcolor[rgb]{0.56,0.35,0.01}{\textbf{\textit{#1}}}}
\newcommand{\AlertTok}[1]{\textcolor[rgb]{0.94,0.16,0.16}{#1}}
\newcommand{\ErrorTok}[1]{\textcolor[rgb]{0.64,0.00,0.00}{\textbf{#1}}}
\newcommand{\NormalTok}[1]{#1}

\newlength{\cslhangindent}
\setlength{\cslhangindent}{1.5em}
\newlength{\csllabelwidth}
\setlength{\csllabelwidth}{3em}
\newenvironment{CSLReferences}[3] % #1 hanging-ident, #2 entry spacing
 {% don't indent paragraphs
  \setlength{\parindent}{0pt}
  % turn on hanging indent if param 1 is 1
  \ifodd #1 \everypar{\setlength{\hangindent}{\cslhangindent}}\ignorespaces\fi
  % set entry spacing
  \ifnum #2 > 0
  \setlength{\parskip}{#2\baselineskip}
  \fi
 }%
 {}
\usepackage{calc} % for \widthof, \maxof
\newcommand{\CSLBlock}[1]{#1\hfill\break}
\newcommand{\CSLLeftMargin}[1]{\parbox[t]{\maxof{\widthof{#1}}{\csllabelwidth}}{#1}}
\newcommand{\CSLRightInline}[1]{\parbox[t]{\linewidth}{#1}}
\newcommand{\CSLIndent}[1]{\hspace{\cslhangindent}#1}\usepackage[table,dvipsnames]{xcolor}

\geometry{includemp,
            letterpaper,
            top=2.4cm,
            bottom=2.4cm,
            left=1.0cm,
            right=1.0cm,
            marginparwidth=5cm,
            marginparsep=1.0cm}

\usepackage[singlelinecheck=off]{caption}

\captionsetup{
  font={small},
  labelfont={bf},
  format=plain,
  labelsep=quad
}

\usepackage{floatrow}

\floatsetup[figure]{margins=hangright,
              facing=no,
              capposition=beside,
              capbesideposition={center,outside},
              floatwidth=\textwidth}

% \floatsetup[table]{margins=hangright,
%              facing=no,
%              capposition=beside,
%              capbesideposition={center,outside},
%              floatwidth=\textwidth}

\pagestyle{plain}

\setcounter{secnumdepth}{5}

\usepackage{titlesec}

\titleformat{\section}[block]
{\normalfont\large\sffamily}
{\thesection}{.5em}{\titlerule\\[.8ex]\bfseries}

\titleformat{\subsection}[runin]
{\normalfont\fontseries{b}\selectfont\filright\sffamily}
{\thesubsection.}{.5em}{}

\titleformat{\subsubsection}[runin]
{\normalfont\itshape\rmfamily\bfseries}{\thesubsubsection}{1em}{}

\fancypagestyle{firstpage}
{
   \fancyhf{}
   \renewcommand{\headrulewidth}{0pt}
   \fancyfoot[R]{\footnotesize\ccby}
   \fancyfoot[L]{\footnotesize\sffamily\today}
}

\fancypagestyle{normal}
{
  \fancyhf{}
  \fancyfoot[R]{\footnotesize\sffamily\thepage\ of \pageref*{LastPage}}
}

\usepackage{tikz}
\begin{document}
\tikz [remember picture, overlay] %
\node [shift={(-0.6in,1.1cm)},scale=0.2,opacity=0.4] at (current page.south east)[anchor=south east]{\includegraphics{logo}};%
\pagestyle{normal}
\thispagestyle{firstpage}

\newcommand{\colorRule}[3][black]{\textcolor[HTML]{#1}{\rule{#2}{#3}}}

\noindent {\LARGE \textbf{\textsf{The biological interpretation of
probabilistic ecological networks}}}

\medskip
\begin{flushleft}
{\small
%
\href{https://orcid.org/0000-0001-9051-0597}{Francis\,Banville}%
%
\,\textsuperscript{1,2,3,‡}, %
\href{https://orcid.org/0000-0001-6067-1349}{Tanya\,Strydom}%
%
\,\textsuperscript{1,3,‡}, %
\href{https://orcid.org/0000-0002-0735-5184}{Timothée\,Poisot}%
%
\,\textsuperscript{1,3}
\vskip 1em
\textsuperscript{1}\,Université de
Montréal; \textsuperscript{2}\,Université de
Sherbrooke; \textsuperscript{3}\,Quebec Centre for Biodiversity
Science\\
\textsuperscript{‡}\,These authors contributed equally to the work\\
\vskip 1em
\textbf{Correspondance to:}\\
Francis Banville --- \texttt{francis.banville@umontreal.ca}\\
}
\end{flushleft}

\vskip 2em
\makebox[0pt][l]{\colorRule[CCCCCC]{2.0\textwidth}{0.5pt}}
\vskip 2em
\noindent

\marginpar{\vskip 1em\flushright
{\small{\bfseries Keywords}:\par
ecological networks\\metaweb\\probabilistic networks\\spatial
scale\\species interactions\\temporal scale\\}
}




        {\bfseries 1:}\,Community ecologists are increasingly shifting
from a binary thinking of ecological networks (e.g., do species
interact?) to a more probabilistic perspective (e.g., how likely are
species to interact?). Assuredly, the benefits of representing
ecological interactions as probabilistic events are numerous, from a
better assessment of the spatial variation of interactions to an
increase capacity to reconstruct networks from sparse data.\\%
        {\bfseries 2:}\,However, probabilities need to be used with
caution when working with interactions and ecological networks. Indeed,
depending on the system at hand and the method used to build
probabilistic networks, probabilities can have different interpretations
that imply different ways to manipulate them. This is rarely discussed
in the literature, thus impeding our ability to use data on
probabilistic interactions appropriately.\\%
        {\bfseries 3:}\,At the core of these differences lie the
distinction between assessing the likelihood that two groups of
individuals \emph{can} interact and the likelihood that they \emph{will}
interact. This impacts the spatial, temporal, and taxonomic scaling of
interaction probabilities, thus further enlightening the need to
properly define them in their ecological context.\\%
        {\bfseries 4:}\,With these challenges in mind, we propose a
general approach to thinking about probabilities in regards to
ecological interactions, and call for better definitions and
conceptualizations of probabilistic ecological networks, both at the
local and regional scales.\\%
    

\vskip 2em
\makebox[0pt][l]{\colorRule[CCCCCC]{2.0\textwidth}{0.5pt}}
\vskip 2em

\hypertarget{introduction}{%
\section{Introduction}\label{introduction}}

\hypertarget{general-background}{%
\subsection{General background}\label{general-background}}

Why is it useful to think about interactions as probabilistic events?

\begin{itemize}
\tightlist
\item
  A biological interaction is probabilistic since two taxa co-occurring
  does not mean they are going to meet (e.g., think of species relative
  abundances). Also, two individuals meeting does not mean that an
  interaction will occur (e.g., a lion crossing paths with a gazelle
  does not mean predation).
\item
  Interaction is contextual on the environment and on the physiological
  state of both species (or individuals).
\item
  Representing ecological interactions as probabilistic events helps us
  predict networks across time and space and assess the spatial
  variability of biological interactions.
\end{itemize}

\textbf{Papers:} Poisot \emph{et al.} (2016)

\hypertarget{problem-and-objectives}{%
\subsection{Problem and objectives}\label{problem-and-objectives}}

Why should we use probabilities with caution when working with
ecological networks? What is the objective of this paper?

\begin{itemize}
\tightlist
\item
  There are different ways to define and interpret interaction
  probabilities depending on the study system and on the method used to
  build probabilistic networks.
\item
  It is important to document and define what we mean by an interaction
  probability because different definitions can have different
  ecological and statistical implications and final interpretation.
\item
  This paper aims to outline some of the challenges in interpreting
  interaction probabilities and propose an approach to thinking about
  them.
\end{itemize}

\hypertarget{definitions-and-interpretations}{%
\section{Definitions and
interpretations}\label{definitions-and-interpretations}}

\hypertarget{overview-of-interaction-probabilities}{%
\subsection{Overview of interaction
probabilities}\label{overview-of-interaction-probabilities}}

How are interaction probabilities defined in the literature? It might
not be as intuitive as one would think.

\begin{itemize}
\tightlist
\item
  There is a big difference in how we interpret the probability that two
  species \emph{can} interact (metaweb) and the probability that they
  \emph{will} interact (realized networks).
\item
  Interaction probabilities can be used to describe Boolean interactions
  (e.g., the probability that two species interact) and weighted
  interactions (e.g., the probability distribution of the amount of
  energy flow between two species).
\item
  In many studies, it is not obvious if their numbers represent
  interaction scores or probabilities (in the latter case, it is rarely
  specified what these probabilities truly represent).
\end{itemize}

\hypertarget{probabilistic-metawebs}{%
\subsection{Probabilistic metawebs}\label{probabilistic-metawebs}}

What does a probability in the context of a metaweb mean?

\begin{itemize}
\tightlist
\item
  It means the probability that two taxa can interact, regardless of
  biological plasticity, environmental variability, or co-occurrence.
\item
  One observation is enough to set this probability to one.
\item
  Can we turn this into a local network realisation that is also
  probabilistic and intuitive?
\end{itemize}

\textbf{Papers:} Strydom \emph{et al.} (2022)

\hypertarget{probabilistic-local-networks}{%
\subsection{Probabilistic local
networks}\label{probabilistic-local-networks}}

What does a probability in the context of a local network mean? A
cautionary tale of how we define probabilities.

\begin{itemize}
\tightlist
\item
  It means the probability that two taxa will interact, regardless of
  biological plasticity and environmental variability.
\item
  What do we mean by two taxa will interact? That at least one
  individual from one group will interact with (e.g., predate) at least
  one other individual from the other group.
\item
  The probability is conditional on the environmental and local
  abundance contexts.
\item
  We should expect a certain number of interactions to occur depending
  on the probability value. This number depends on the number of trials,
  which is defined according to the ecological context (e.g.,
  environmental conditions, scale) in which probabilities were
  estimated. This is in contrast with probabilities in metawebs.
\end{itemize}

\hypertarget{scaling}{%
\section{Scaling}\label{scaling}}

\hypertarget{spatial-and-temporal-scales}{%
\subsection{Spatial and temporal
scales}\label{spatial-and-temporal-scales}}

How do interaction probabilities scale spatially and temporally?

\begin{itemize}
\tightlist
\item
  Why do probabilistic local networks scale with area and time but not
  probabilistic metawebs?
\item
  Network area relationships (e.g., Ontario lakes and Alaska datasets)
\item
  In metawebs, interaction probabilities do not scale with space and
  time.
\item
  In local networks, interaction probabilities scale with space and
  time. There are more opportunities of interactions (e.g., more
  environmental conditions) in a larger area and longer time period.
\item
  Regional networks can inform local networks but can local networks
  inform regional networks?
\end{itemize}

\textbf{Figure:} Empirical example of the association between the number
of interactions in realized local networks and the number of
interactions in the corresponding species subnetworks of regional
networks. We should expect the interaction to be linear below the 1:1
line.

\textbf{Papers:} there might be something in these McLeod \emph{et al.}
(2020); McLeod \emph{et al.} (2021); Wood \emph{et al.} (2015)

\hypertarget{taxonomic-scale}{%
\subsection{Taxonomic scale}\label{taxonomic-scale}}

How do interaction probabilities scale taxonomically?

\begin{itemize}
\tightlist
\item
  Difference of interpretation between probabilities in networks at the
  individual level and at any higher taxonomic level.
\item
  How does the scaling up of the nodes from an individual to population
  to taxonomic group group change our interpretation of interaction
  probabilities? How does the aggregation change the interpretation?
\item
  How is it analogous to spatial and temporal scaling? Basically, all
  kinds of scaling are just different ways to aggregate individuals or
  nodes.
\end{itemize}

\textbf{Figure:} Conceptual figure of how a scale up of the nodes from
an individual to a population to any taxonomic group change our
interpretation of the probability of interaction.

\hypertarget{concluding-remarks}{%
\section{Concluding remarks}\label{concluding-remarks}}

Here we present some advice moving forward.

\begin{itemize}
\tightlist
\item
  Non-ecological Networks: What can we learn from other systems/fields
  (e.g., social networks, probabilistic graph theory)?
\item
  What even are the probabilities? What is the probability that we will
  ever know the answer to that?
\item
  Be careful how we define probabilities. Be sure to be explicit about
  these things. Be sure to specify the type of interaction, the spatial,
  temporal, and taxonomic scale when presenting new data on interaction
  probabilities. We need better metadata documentation.
\item
  Be careful to use and manipulate interaction probabilities properly
  depending on how they were defined. Different interpretations imply
  different scaling, and thus different ways to manipulate these
  numbers.
\item
  Maybe mention thinking about a workflow to predict probabilistic local
  networks from probabilistic metawebs.
\end{itemize}

\hypertarget{references}{%
\section*{References}\label{references}}
\addcontentsline{toc}{section}{References}

\hypertarget{refs}{}
\begin{CSLReferences}{1}{0}
\leavevmode\hypertarget{ref-McLeod2021SamAsy}{}%
McLeod, A., Leroux, S.J., Gravel, D., Chu, C., Cirtwill, A.R., Fortin,
M.-J., \emph{et al.} (2021). Sampling and asymptotic network properties
of spatial multi-trophic networks. \emph{Oikos}, n/a.

\leavevmode\hypertarget{ref-McLeod2020EffSpe}{}%
McLeod, A.M., Leroux, S.J. \& Chu, C. (2020). Effects of species traits,
motif profiles, and environment on spatial variation in multi-trophic
antagonistic networks. \emph{Ecosphere}, 11, e03018.

\leavevmode\hypertarget{ref-Poisot2016StrProa}{}%
Poisot, T., Cirtwill, A.R., Cazelles, K., Gravel, D., Fortin, M.-J. \&
Stouffer, D.B. (2016). The structure of probabilistic networks.
\emph{Methods in Ecology and Evolution}, 7, 303--312.

\leavevmode\hypertarget{ref-Strydom2022FooWeb}{}%
Strydom, T., Bouskila, S., Banville, F., Barros, C., Caron, D., Farrell,
M.J., \emph{et al.} (2022). Food web reconstruction through phylogenetic
transfer of low-rank network representation. \emph{Methods in Ecology
and Evolution}, n/a.

\leavevmode\hypertarget{ref-Wood2015EffSpa}{}%
Wood, S.A., Russell, R., Hanson, D., Williams, R.J. \& Dunne, J.A.
(2015). Effects of spatial scale of sampling on food web structure.
\emph{Ecology and Evolution}, 5, 3769--3782.

\end{CSLReferences}

\end{document}
